% GENERATED FILE. Edit canonical lessons, then run npm run generate:books.
% canonical-source-hash: fnv1a64:d9f124a47a0a770d
% canonical-lessons: UR-C115-adrak, UR-C115-gobhi, UR-C115-matar, UR-C115-khira, UR-C115-tarbuz

\chapter{Ginger, Cauliflower, Peas, Cucumber, Watermelon}
\label{ch:ur-ginger-cauliflower-peas-cucumber-watermelon}
\hlchaptermodality{115}

\begin{chapteropening}
\textbf{By the end of this chapter:} \emph{I can say ginger, a cauliflower, peas, a cucumber and a watermelon.}

Complete the last lesson of chapter 115: I can say ginger, a cauliflower, peas, a cucumber and a watermelon.
\end{chapteropening}

\section[\texorpdfstring{\ur{ادرک} (adrak)}{ادرک (adrak)}]{\ur{ادرک} (adrak) — ginger}

\label{lesson:UR-C115-adrak}



\begin{quote}
Before the new one: say the Urdu for an onion, then the Urdu for garlic.
\end{quote}

\subsection*{You'll want to know: \ur{ادرک}}
\textbf{\ur{ادرک}} — \emph{adrak} — \textquotedblleft{}ginger\textquotedblright{}.

\begin{cousinweb}
\textbf{\ur{ادرک}} (\emph{adrak}) goes back to Sanskrit \emph{ārdraka}, \textquotedblleft{}fresh ginger\textquotedblright{} (CDIAL 1341). CDIAL lists the inherited Hindi form as \emph{ādā}; \emph{adrak} keeps more of the Sanskrit shape, so it was probably taken over later and lightly reshaped, a semi-tatsama.
\end{cousinweb}

\begin{grammarlens}[title={What you've built}]
One more word for this chapter.
\end{grammarlens}

\subsection*{Your turn}
\begin{itemize}
\raggedright
  \item \emph{Say it:} \emph{adrak}
  \item \emph{Say it:} \emph{adrak}, once more
  \item \emph{Say it:} say \emph{adrak}
  \item \emph{Recall:} say the Urdu for an onion, then the Urdu for garlic, then say \emph{adrak} again
\end{itemize}

\begin{tcolorbox}[breakable,colback=teal!4,colframe=teal!35!black,title={Before you move on}]
What does \ur{ادرک} mean? (\textquotedblleft{}Ginger\textquotedblright{}.) Say it once more.
\end{tcolorbox}
\section[\texorpdfstring{\ur{گوبھی} (gobhī)}{گوبھی (gobhī)}]{\ur{گوبھی} (gobhī) — a cauliflower}

\label{lesson:UR-C115-gobhi}



\begin{quote}
Before the new one: say the Urdu for garlic, then the Urdu for ginger.
\end{quote}

\subsection*{You'll want to know: \ur{گوبھی}}
\textbf{\ur{گوبھی}} — \emph{gobhī} — \textquotedblleft{}a cauliflower\textquotedblright{}.

\begin{cousinweb}
\textbf{\ur{گوبھی}} (\emph{gobhī}) is a tadbhava, a word inherited from Sanskrit and reshaped by sound change. It comes from \emph{gojihvikā}, \textquotedblleft{}a kind of plant\textquotedblright{} (CDIAL 4270). The Sanskrit name means \textquotedblleft{}cow's tongue\textquotedblright{}; CDIAL glosses the Hindi word \textquotedblleft{}cabbage\textquotedblright{}.
\end{cousinweb}

\begin{grammarlens}[title={What you've built}]
One more word for this chapter.
\end{grammarlens}

\subsection*{Your turn}
\begin{itemize}
\raggedright
  \item \emph{Say it:} \emph{gobhī}
  \item \emph{Say it:} \emph{gobhī}, once more
  \item \emph{Say it:} say \emph{gobhī}
  \item \emph{Recall:} say the Urdu for garlic, then the Urdu for ginger, then say \emph{gobhī} again
\end{itemize}

\begin{tcolorbox}[breakable,colback=teal!4,colframe=teal!35!black,title={Before you move on}]
What does \ur{گوبھی} mean? (\textquotedblleft{}A cauliflower\textquotedblright{}.) Say it once more.
\end{tcolorbox}
\section[\texorpdfstring{\ur{مٹر} (maṭar)}{مٹر (maṭar)}]{\ur{مٹر} (maṭar) — peas}

\label{lesson:UR-C115-matar}



\begin{quote}
Before the new one: say the Urdu for ginger, then the Urdu for a cauliflower.
\end{quote}

\subsection*{You'll want to know: \ur{مٹر}}
\textbf{\ur{مٹر}} — \emph{maṭar} — \textquotedblleft{}peas\textquotedblright{}.

\begin{cousinweb}
\textbf{\ur{مٹر}} (\emph{maṭar}) is a tadbhava, a word inherited and reshaped by sound change. It comes from \emph{*maṭṭara}, \textquotedblleft{}pea\textquotedblright{}, a form CDIAL reconstructs from its descendants (CDIAL 9724).
\end{cousinweb}

\begin{grammarlens}[title={What you've built}]
One more word for this chapter.
\end{grammarlens}

\subsection*{Your turn}
\begin{itemize}
\raggedright
  \item \emph{Say it:} \emph{maṭar}
  \item \emph{Say it:} \emph{maṭar}, once more
  \item \emph{Say it:} say \emph{maṭar}
  \item \emph{Recall:} say the Urdu for ginger, then the Urdu for a cauliflower, then say \emph{maṭar} again
\end{itemize}

\begin{tcolorbox}[breakable,colback=teal!4,colframe=teal!35!black,title={Before you move on}]
What does \ur{مٹر} mean? (\textquotedblleft{}Peas\textquotedblright{}.) Say it once more.
\end{tcolorbox}
\section[\texorpdfstring{\ur{کھیرا} (khīrā)}{کھیرا (khīrā)}]{\ur{کھیرا} (khīrā) — a cucumber}

\label{lesson:UR-C115-khira}



\begin{quote}
Before the new one: say the Urdu for a cauliflower, then the Urdu for peas.
\end{quote}

\subsection*{You'll want to know: \ur{کھیرا}}
\textbf{\ur{کھیرا}} — \emph{khīrā} — \textquotedblleft{}a cucumber\textquotedblright{}.

\begin{cousinweb}
\textbf{\ur{کھیرا}} (\emph{khīrā}) is a tadbhava, a word inherited from Sanskrit and reshaped by sound change. It comes from \emph{kṣīraka}, \textquotedblleft{}a fragrant plant\textquotedblright{} (CDIAL 3697). It is built on \emph{kṣīra}, \textquotedblleft{}milk\textquotedblright{}; the Sanskrit word is known only from the old word-lists.
\end{cousinweb}

\begin{grammarlens}[title={What you've built}]
One more word for this chapter.
\end{grammarlens}

\subsection*{Your turn}
\begin{itemize}
\raggedright
  \item \emph{Say it:} \emph{khīrā}
  \item \emph{Say it:} \emph{khīrā}, once more
  \item \emph{Say it:} say \emph{khīrā}
  \item \emph{Recall:} say the Urdu for a cauliflower, then the Urdu for peas, then say \emph{khīrā} again
\end{itemize}

\begin{tcolorbox}[breakable,colback=teal!4,colframe=teal!35!black,title={Before you move on}]
What does \ur{کھیرا} mean? (\textquotedblleft{}A cucumber\textquotedblright{}.) Say it once more.
\end{tcolorbox}
\section[\texorpdfstring{\ur{تربوز} (tarbūz)}{تربوز (tarbūz)}]{\ur{تربوز} (tarbūz) — a watermelon}

\label{lesson:UR-C115-tarbuz}



\begin{quote}
Before the new one: say the Urdu for peas, then the Urdu for a cucumber.
\end{quote}

\subsection*{You'll want to know: \ur{تربوز}}
\textbf{\ur{تربوز}} — \emph{tarbūz} — \textquotedblleft{}a watermelon\textquotedblright{}.

\begin{cousinweb}
\textbf{\ur{تربوز}} (\emph{tarbūz}) is a loan, usually traced to Persian \emph{tarbuz}, \textquotedblleft{}watermelon\textquotedblright{}.
\end{cousinweb}

\begin{grammarlens}[title={What you've built}]
One more word for this chapter.
\end{grammarlens}

\subsection*{Your turn}
\begin{itemize}
\raggedright
  \item \emph{Say it:} \emph{tarbūz}
  \item \emph{Say it:} \emph{tarbūz}, once more
  \item \emph{Say it:} say \emph{tarbūz}
  \item \emph{Recall:} say the Urdu for peas, then the Urdu for a cucumber, then say \emph{tarbūz} again
\end{itemize}

\begin{tcolorbox}[breakable,colback=teal!4,colframe=teal!35!black,title={Before you move on}]
What does \ur{تربوز} mean? (\textquotedblleft{}A watermelon\textquotedblright{}.) Say it once more.
\end{tcolorbox}
